\section{Corrections after the second independent review}\label{sec:reviewtwocorrections}
This section corrects claims in the 30 September manuscript. The unaltered
second review recommends rejection at confidence 4/5. The new mathematical
observations and ongoing experiments below have not received a further full
review. No acceptance or useful end-to-end uncertainty result is asserted.

\subsection{Matched continuous Gaussian inference}
Write $G=\A\A^*$, $a=\A\ell$, and $L=\norm{\ell}^2$ for a fixed, whitened
continuous design. A generalized white Gaussian prior of directional scale
$\tau$ has weights $w_\tau=(G+\tau^{-2}I)^{-1}a$ and functional variance
\begin{align}
v_\tau&=\tau^2(L-a^\top w_\tau)\nonumber\\
&=\norm{w_\tau}^2+\tau^2\norm{\ell-\A^*w_\tau}^2.
\end{align}
White covariance on infinite-dimensional $L^2$ is not trace class: this is an
isonormal process on test functions, not an $L^2$-valued density with finite
energy. The finite joint distribution of observations and target is nevertheless
well defined. The construction is continuous ridge inference; neither that
identity nor the resolving-kernel viewpoint is new.

\begin{proposition}[A fixed-bias normal envelope]
For $z\ge\sqrt{3}$ and $t\ge0$,
$\Pr(|N(t,1)|\le z\sqrt{1+t^2})\ge2\Phi(z)-1$.
Thus $z_{1-\alpha/2}\sqrt{s^2+b^2}$ covers a Gaussian error with variance
$s^2$ and any fixed absolute bias at most $b$, whenever
$\alpha\le2\Phi(-\sqrt{3})$.
\end{proposition}
\begin{proof}
Set $q=z\sqrt{1+t^2}$, $F(t)=\Phi(q-t)+\Phi(q+t)-1$ and
$c=zt/\sqrt{1+t^2}$. Then
\[
F'(t)=\phi(q-t)[(c-1)+(c+1)e^{-2qt}].
\]
For $c\ge1$ this is nonnegative. For $c<1$ this is equivalent to
$\operatorname{arctanh}(c)\ge qt=c/(1-c^2/z^2)$. The absolutely convergent
series have coefficients $1/(2j+1)$ and $z^{-2j}$, respectively.
Since $z^{2j}\ge3^j\ge2j+1$, the inequality follows termwise.
Therefore $F(t)\ge F(0)=2\Phi(z)-1$. Zero noise follows directly.
\end{proof}
At $\tau=B$, the usual Gaussian prior interval therefore shares the uniform
fixed-pose guarantee over the $B$-ball. The folded-normal value can shorten it,
and different weights can change efficiency, but uniform fixed-pose coverage
is not a distinction from this matched baseline. At $\tau=B/2$ the implication
need not hold. Our earlier finite-grid controls had a different operator and
smaller directional prior scales. Their behavior cannot establish superiority
over continuous Gaussian inference. The 48-fit matched comparison is running;
its first higher-band solve hits its iteration limit and is retained as
nonconverged, not an exact posterior calculation.

\subsection{Noise-dependent nuisance errors}
Theorem~\ref{thm:coverage} requires nuisance errors fixed independently of the
inference noise. Even with fixed weights, a bounded bias chosen as
$b\,\operatorname{sign}(Z)$ gives absolute error $s|Z|+b$.
The folded-normal interval need not cover it at the advertised level.
The sum width $b+z_{1-\alpha/2}s$ handles this bounded adaptive bias if the
noise projection remains fixed Gaussian. It does not repair weights and
poses estimated from the same noisy image. An independent inference image
makes that projection conditionally Gaussian, but does not by itself make
alignment errors centered conditional on the fitted design.

\subsection{Prospective local-alignment study}
A separate calibration uses 128 noisy replicate datasets per acquisition
geometry and both an oracle-reference and an independent-pilot local alignment
template. The frozen follow-up uses 200 additional datasets per geometry,
recomputing poses and weights, with same-image and independent-image inference
controls. It compares fixed-pose, deterministic-pose, centered mixed-pose and
matched continuous Gaussian intervals on two broad targets. All failed fits,
raw intervals and no-data fallbacks are retained. This is local alignment
against fixed templates, not a complete ab initio reconstruction pipeline.
No completed coverage result from this study is included in this revision.
